\documentclass[11pt,a4paper]{article}
\usepackage[utf8]{inputenc}
\usepackage[T1]{fontenc}
\usepackage{lmodern,amsmath,amssymb,amsthm,mathtools,array,booktabs,enumitem,microtype}
\usepackage{natbib}
\usepackage[margin=1in]{geometry}
\usepackage{xr-hyper}
\usepackage{hyperref}
\externaldocument[GE-]{01_law_of_graphic_equalisation_theory_referenced}
\externaldocument[MS-]{01b_graphic_equalisation_mind_and_self_grounding}
\newtheorem{definition}{Definition}[section]
\newtheorem{proposition}[definition]{Proposition}
\newtheorem{corollary}[definition]{Corollary}
\theoremstyle{remark}\newtheorem{remark}[definition]{Remark}
\title{The Law of Graphic Equalisation --- Communication Grounding\\\large Difference Across Frames, Kinetic Contact, Records, and Network Frames}
\author{Matthew Parslow\\Independent Researcher}
\date{\textbf{Draft \texorpdfstring{\textemdash}{---}\ \PaperDraft}\\[2pt]\PaperDisplayDate \texorpdfstring{\textbullet}{\textperiodcentered}\ document revision \PaperRevision}
\usepackage{fancyhdr}
\pagestyle{fancy}
\fancyhf{}
\fancyfoot[C]{\thepage}
\fancyfoot[R]{\scriptsize\textsc{\PaperDraft\,r\PaperRevision}}
\renewcommand{\headrulewidth}{0pt}
\fancypagestyle{plain}{\fancyhf{}\fancyfoot[C]{\thepage}%
  \fancyfoot[R]{\scriptsize\textsc{\PaperDraft\,r\PaperRevision}}\renewcommand{\headrulewidth}{0pt}}

% GENERATED by tools/build_paper_metadata.py from paper_metadata.json.
% DO NOT EDIT: edit paper_metadata.json, then regenerate this file.
\providecommand{\PaperDraft}{draft-2026-09-06.1}
\providecommand{\PaperDisplayDate}{6 September 2026}
\expandafter\def\csname PaperRevision@00_is\endcsname{129}
\expandafter\def\csname PaperRevision@01_law_of_graphic_equalisation_theory_referenced\endcsname{294}
\expandafter\def\csname PaperRevision@01a_graphic_equalisation_physics_grounding\endcsname{47}
\expandafter\def\csname PaperRevision@01b_graphic_equalisation_mind_and_self_grounding\endcsname{47}
\expandafter\def\csname PaperRevision@01c_graphic_equalisation_cosmology_grounding\endcsname{11}
\expandafter\def\csname PaperRevision@01d_graphic_equalisation_born_recovery\endcsname{6}
\expandafter\def\csname PaperRevision@01e_graphic_equalisation_communication_grounding\endcsname{6}
\expandafter\def\csname PaperRevision@01f_graphic_equalisation_qualia_and_consciousness_grounding\endcsname{3}
\expandafter\def\csname PaperRevision@01g_graphic_equalisation_total_perspective_vortex_grounding\endcsname{4}
\expandafter\def\csname PaperRevision@01h_graphic_equalisation_information_gravity_quantum_reconciliation\endcsname{5}
\expandafter\def\csname PaperRevision@01i_graphic_equalisation_passive_continuation_momentum_gravity\endcsname{12}
\expandafter\def\csname PaperRevision@01s_graphic_equalisation_statement\endcsname{294}
\expandafter\def\csname PaperRevision@02_seldon_manifolds_theory_referenced\endcsname{31}
\expandafter\def\csname PaperRevision@02a_seldon_manifolds_psychological_grounding\endcsname{15}
\expandafter\def\csname PaperRevision@02b_seldon_manifolds_social_and_psychological_application\endcsname{7}
\expandafter\def\csname PaperRevision@02c_seldon_manifolds_decision_making\endcsname{11}
\expandafter\def\csname PaperRevision@03_law_of_narrative_momentum_theory_referenced\endcsname{32}
\expandafter\def\csname PaperRevision@04_narrativium_theory_referenced\endcsname{23}
\expandafter\def\csname PaperRevision@05_frames_relation_and_the_empty_limit\endcsname{11}
\expandafter\def\csname PaperRevision@06_assessing_an_account_of_this_kind\endcsname{29}
\expandafter\def\csname PaperRevision@07_causal_attitude_networks_on_the_substrate\endcsname{4}
\expandafter\def\csname PaperRevision@08_certified_answers_as_an_audit\endcsname{4}
\expandafter\def\csname PaperRevision@09_adaptability\endcsname{94}
\expandafter\def\csname PaperRevision@09a_adaptability_implementation\endcsname{7}
\expandafter\def\csname PaperRevision@10_applied_graphic_equalisation\endcsname{7}
\expandafter\def\csname PaperRevision@10a_red_queen_race_civilisational_rate_of_change_llms_biological_timescales_anthropocene_transition\endcsname{13}
\expandafter\def\csname PaperRevision@11_relational_information_ethics\endcsname{16}
\expandafter\def\csname PaperRevision@A1_constructive_graphic_equalisation_implementation_referenced\endcsname{8}
\expandafter\def\csname PaperGenerated@01s_graphic_equalisation_statement\endcsname{1}
\edef\PaperMetadataStem{\jobname}
\expandafter\ifx\csname PaperRevision@\PaperMetadataStem\endcsname\relax
  \PackageError{paper_metadata}{No metadata for \PaperMetadataStem}{Add the paper stem to paper_metadata.json.}
  \def\PaperRevision{?}
\else
  \expandafter\let\expandafter\PaperRevision\csname PaperRevision@\PaperMetadataStem\endcsname
\fi
\newcommand{\PaperDateLine}{\textbf{Draft \texorpdfstring{\textemdash}{---}\ \PaperDraft}\\[2pt]\PaperDisplayDate\ \texorpdfstring{\textbullet}{\textperiodcentered}\ document revision \PaperRevision}
\newcommand{\PaperStatementDateLine}{\textbf{Draft \texorpdfstring{\textemdash}{---}\ \PaperDraft}\\[2pt]\PaperDisplayDate\ \texorpdfstring{\textbullet}{\textperiodcentered}\ statement}
\expandafter\ifx\csname PaperGenerated@\PaperMetadataStem\endcsname\relax
  \newcommand{\PaperFooterStamp}{\textsc{\PaperDraft\,r\PaperRevision}}
\else
  \newcommand{\PaperFooterStamp}{\textsc{\PaperDraft\ / statement / from long r\PaperRevision}}
\fi
\newcommand{\PaperNoticeStamp}{\textbf{DRAFT --- \PaperDraft\ r\PaperRevision} (\PaperDisplayDate).}


\begin{document}
\maketitle
\begin{center}
\fbox{\parbox{0.92\textwidth}{\small
\textbf{DRAFT --- \PaperDraft r\PaperRevision} (\PaperDisplayDate).
This document is under active revision and is circulated for comment, not as a
finished result. It is the \emph{01e} grounding companion to \emph{The Law of
Graphic Equalisation --- Theory} (Graphic Equalisation), and inherits that
paper's definitions, epistemic boundary, and level-of-claims. It does not add a
primitive to the substrate, assign a communication ontology over it, or alter
Graphic Equalisation's core claims. It develops a cross-frame grounding: all
causal influence is communication at some grain, while communication may be
held, offered, or carried without producing a relevant effect in a particular
receiving frame. The sections on falsifiability identify promoted commitments,
not a proof of the substrate.}}
\end{center}
\vspace{1em}

\begin{abstract}
Communication is usually split too early: physical influence is put on one side,
signals and language on another, and the latter is made to depend on a human
speaker, a shared intention, or an intrinsically meaningful symbol. This paper
uses Graphic Equalisation to refuse that split. A tap on the shoulder, a punch,
a molecular binding event, a written record, an Ethernet transmission, and a
word-token are all differences carried, exposed, or held across frame relations.
They differ in magnitude, grain, persistence, grammar, and what another frame
can make of them, not in one being communication and another being merely
physical cause. Causal influence is communication realised as a transition in a
receiving frame. Conversely, a record, utterance, or offered packet can
communicate substantial structure while producing no relevant uptake or action
in a particular frame.

The paper's governing constraint is economic: communication pays where it
changes frames. Carriage, reception, and repair are distinct writes, charged to
the frame relation in which each occurs; a later description or a containing
frame may not silently provide distinctions free to the frame being described.
Network link-layer frames provide the first central engineering grounding.
Their delimiters, length or flag grammar, headers, payloads, integrity checks,
discard paths, buffers, acknowledgement states, and encapsulation show literal,
layer-relative frames in operation. They also impose limits: a checksum is not
meaning, a frame is not a network-layer datagram, and endpoint correctness cannot
be settled by a lower layer. Human lexical retrieval and tip-of-the-tongue are
then a bounded receiver-side case of a held shape whose lexical edge is
unavailable, rather than the definition of communication.
\end{abstract}

\section{Scope and claim}\label{sec:scope}

This is a grounding paper, not a new substrate and not a psychology of
communication. It inherits Graphic Equalisation's minimal graph $G=(g,r,\Gamma)$,
its frame-relative grounding (Axiom~\ref{GE-ax:no-view}), its account of
admissibility and residual, and its write-economics. A frame is real graph data
that runs (Graphic Equalisation's Definition~\ref{GE-def:frame}); a frame's
aperture and grammar determine what it can take in and what it can do with what
it takes in. The present task is to show that cross-frame communication can be
stated entirely in those inherited terms. It does not posit a semantic substance,
an intrinsic message, a universal codebook, a privileged observer, or a distinct
communication primitive beneath relation.

The strongest claim is structural. A difference can be carried from one frame
such that another frame can hold, recover, classify, or act on it, and this is a
relation between frames rather than a property a token possesses alone. The
weaker and empirical claims concern which grammars permit which recovery, what
has to be retained for repair, what costs are paid at which crossings, and how
particular physical, biological, computational, and linguistic systems realise
the relation. The evidence for a composition claim is recurrence of the same
shape under independent domain descriptions; that recurrence is neither proof of
the substrate nor proof that this paper has found an exhaustive taxonomy.

The paper also corrects a familiar but unhelpful opposition. There is no
substrate-level divide between causal influence and communication. A punch is
not physical influence with communication absent, while a tap on the shoulder is
communication because a human can interpret it. Both are kinetic communication:
differences in pressure, momentum, direction, location, duration, and consequence
cross a bodily boundary. They differ in magnitude and in what further frames can
do with the event, not in kind. A tap may be low-energy and support a rich social
reconstruction; a punch may be high-energy and force a large state transition.
Neither has a meaning independent of every frame, and neither requires an
intentional source in order to be communication.

\begin{definition}[Communication]\label{def:communication}
Communication is a distinction carried, exposed, or held across a frame relation.
A \emph{causal influence} is communication realised as a transition in a
receiving frame. An \emph{operative reception} is the narrower case in which the
receiver's rules make the difference available for retention, recovery,
classification, correction, or action.
\end{definition}

The three verbs in the definition are deliberately different. A carried
difference traverses a relation as a realised event. An exposed difference is
available at a boundary under some possible crossing, whether or not this frame
currently takes it in. A held difference persists as graph structure or a record
across the succession relevant to a later crossing. A letter in an unopened
archive, an unqueried database row, and a packet retained in a buffer can all
communicate in the latter senses without producing the action a writer, sender,
or reader might have wanted. That lack of achievement is not a retroactive
annulment of the distinction carried. It only says that an intended or relevant
operative reception did not occur in the frame under examination.

\begin{remark}[Communication is broad; uptake is not]\label{rem:breadth}
The definition makes all causal influence communication at the grain at which it
is a carried difference. It does not make every communication an immediately
realised causal influence in every possible receiver. This is not wordplay. It
keeps separate an inscription and its reading, a protocol offer and its
acceptance, a stored program and its run, a warning and the action it fails to
produce. The distinction is also what prevents an external observer from
crediting a receiver with a difference merely because some containing frame can
read it.
\end{remark}

\section{The economic constraint: payment at a frame change}\label{sec:payment}

Graphic Equalisation's well-formedness condition is that a graph pays for nothing
that does no work for it (Definition~\ref{GE-def:well-formed}); its economics
locates cost in writes rather than in a frame's mere standing structure
(Proposition~\ref{GE-prop:economics}). Communication gives that economy a sharp
location. A difference is not made free to another frame because an external
account can describe source and receiver together. Where it becomes a difference
for that other frame, something changes, and that crossing is where the account
must look for the payment.

\begin{proposition}[Communication pays where it changes frames]\label{prop:pays-at-crossing}
Where a communication becomes a realised, receiver-side difference, the
transition is a write in the receiving frame, in the containing frame that
sustains the crossing, or in both as their rules specify. A later interpretation,
an external description, or a downstream compression may describe the result but
may not silently supply that receiver-side distinction without such a change.
\end{proposition}

This is a promoted application of the inherited economics, not a claim that every
case has one fixed physical price. The payment may be momentum transfer,
deformation, receptor occupation, voltage change, bandwidth, buffer capacity,
parsing work, retained state, a comparison, a repair write, or a rule firing. The
kind and magnitude are contingent rule content. What is ruled out is only the
level error: a receiving frame being credited with a distinction for which no
transition has occurred in it or in the relation that supplies it.

Three payments should be kept distinct.
\begin{enumerate}[leftmargin=*]
\item \emph{Carriage payment} is what makes a difference traverse or remain exposed
at a boundary: pressure through contact, a voltage transition, a diffused molecule,
an ink trace, or a radio signal.
\item \emph{Reception payment} is what makes the difference available at the
receiver's own grain: a receptor state, a membrane potential, a buffer entry, a
parser state, a stored record, or an installed lexical edge.
\item \emph{Repair or reconstruction payment} is what compares alternatives,
registers a residual, corrects an error, requests a retransmission, or recovers a
missing edge from a held shape.
\end{enumerate}

The distinctions are temporal as well as structural. An unread record has paid
for inscription and retention in the holder; it has not yet charged a later
reader for recovery. In the inherited state language this is a realised value that
is not held by the reading frame---undisclosed to it rather than absent or
resolved for it (Definition~\ref{GE-def:states}). A protocol sender may pay to construct and emit a frame; a
receiver pays differently to validate and store it; a retransmission pays only
where an error condition later causes a rule to fire. A record can therefore
communicate much while achieving nothing relevant in the frame that has not read
it. Conversely, the fact that a recipient reacts does not mean it has recovered
every distinction the carrier could make available at another grain.

\begin{remark}[Tap and punch]\label{rem:tap-punch}
A tap and a punch make the rule concrete. Each pays at bodily contact: a kinetic
difference crosses a boundary and changes some local configuration. The punch
usually pays more heavily in momentum transfer, deformation, injury, and later
repair; the tap may pay little kinetically while being rich under a learned social
grammar. Those are differences of scale and of receiver-side structure. It is an
error to call the punch ``mere causation'' and reserve communication for the tap,
or to pretend that the tap's social reading was already present in the force alone.
\end{remark}

\section{The minimal communication relation}\label{sec:minimal-relation}

The proposed construction introduces no object beyond the substrate. It gives
names to roles already supplied by graph relations and rules, because those roles
must be separated if a cross-frame claim is to remain level-correct.

\begin{center}\small
\begin{tabular}{@{}>{\raggedright\arraybackslash}p{0.23\textwidth}>{\raggedright\arraybackslash}p{0.67\textwidth}@{}}
\toprule
\textbf{Communication role} & \textbf{Inherited substrate reading}\\
\midrule
Source-relative distinction & A difference in graph structure, state, or admissible continuation at a named grain.\\
Carrier or record & The realised transition or held graph structure through which that difference is exposed across a relation.\\
Boundary & The relation between frames, whose crossings induce the relevant partial order and whose condition is known only by crossing (Proposition~\ref{GE-prop:boundary-by-crossing}).\\
Receiver grammar & The receiver's $\Gamma$, including any installed recognition, decoding, selection, or response rules.\\
Aperture & The receiver-relative restriction on what can enter and what can be expressed or acted upon.\\
Available distinction & A difference the receiver can retain, reconstruct, discriminate, or condition a rule upon at its own grain.\\
Residual & What remains when an offered or attempted transformation cannot be admitted, recovered, or made coherent under that grammar.\\
Audit frame & A specified frame that can represent the relevant source and receiver relations; it is not a view from nowhere (Axiom~\ref{GE-ax:no-view}).\\
\bottomrule
\end{tabular}
\end{center}

\begin{remark}[These are positional roles, not new graph types]\label{rem:positional-roles}
Source, carrier, boundary, and receiver are positional roles that relations and
adjacent frames occupy at a specified grain; they are not new graph types or a
communication primitive added beneath relation. As in Graphic Equalisation,
naming two ends of a relation differently marks the asymmetry between them, not
two kinds of thing (Corollary~\ref{GE-cor:frame-dimension}). A frame is a source
under one crossing and a receiver under another, and the same edge is carrier to
one grain and boundary to another. These roles are made authoritative in the
stratigraphy of Section~\ref{sec:stratigraphy}: source, sender, receiver,
message, purpose, and intent are audit-band attributions and need not be
represented by any constituent frame of the run
(Remark~\ref{rem:attribution-not-participation}).
\end{remark}

\begin{definition}[Receiver-relative availability]\label{def:availability}
Let $S$ and $R$ be named frames related through a carrier or record $c$. A
source-side distinction is \emph{available at $R$} when $R$'s aperture and grammar
make some receiver-side distinction conditional on $c$. It is \emph{operative at
$R$} when that availability participates in a later rule, retained state, or
realised action. Availability and operation are relative to $R$ and to the grain
specified for it.
\end{definition}

The definition does not say that $c$ is an intrinsically meaningful object, that
$S$ intends anything, or that $R$ reconstructs a source state exactly. It says
only that a relation exists between a specified difference and a receiver-side
one. An audit frame can assess a claimed correspondence only by holding the
relevant relation itself; if it cannot, it has evidence local to one side rather
than a settled mapping. This is the relativity of grounding applied to
communication, not a sceptical denial that messages work.

The same carrier can accordingly do different work at different grains. A printed
sequence is graph structure to a camera sensor, an orthographic token to a reader,
a byte stream to an optical-recognition system, and a decorative mark to a frame
without the relevant aperture. No one of those readings supersedes the others.
They are different communications because different receiver grammars pay for and
make use of different distinctions. The receiving frame's grammar is not a
private mental language: it can be a receptor geometry, a circuit, a protocol
parser, a file format, a learned practice, or a rule in a containing system.

\begin{remark}[No universal codebook]\label{rem:no-universal-codebook}
A shared explicit codebook is neither the definition of communication nor a free
precondition of it. Some receiver grammar must be realised somewhere if a
specified receiver is to discriminate or recover at a specified grain, but it may
be inherited, evolved, learned, installed, negotiated, or distributed across a
larger frame. The paper therefore asks where the grammar is paid for and held; it
does not assume that two communicating frames began with a jointly authored
manual. The labels, tokens, and polarities a grammar uses are supplied by the
frame that reads them rather than being intrinsic to a carrier: Graphic
Equalisation already forces plurality of types without making any type privileged
or self-interpreting (Remark~\ref{GE-rem:types-forced}). This paper accordingly
adds receiver-side economics---where a distinction is paid for, held, and
recovered---rather than a new derivation of the substrate or of its type
plurality. A common grammar or codebook is a frame-relation fact where it is
realised and a social-convention attribution in the audit band
(Section~\ref{sec:stratigraphy}); neither reading makes it a free or universal
precondition of communication.
\end{remark}

\section{The communication stratigraphy: one relation across grains}\label{sec:stratigraphy}

The previous section named the positional roles of a single cross-frame relation.
This section arranges those roles into an explicit hierarchy, or
\emph{stratigraphy}, so that the level at which each concept is a fact can be read
off directly. There is one communication relation; what changes across the
hierarchy is not the relation but the \emph{grain} at which a distinction is a
fact and the \emph{frame} that must hold it. Three bands are distinguished, and
the discipline of the whole section is that a fact stated in a higher band may
never be silently charged to a lower one.

\begin{itemize}[leftmargin=*]
\item The \textbf{substrate band} is where rules run and deformations occur. Its
facts are graph facts: a rule, a deformation, a distinction, a grain, a causal
influence. Nothing in the higher bands adds an operator here; they only describe,
relate, or attribute.
\item The \textbf{frame-relation band} is where a difference is carried, exposed,
or held across a boundary, and where a receiver's aperture and grammar make some
of it available or operative. Its facts are relational and receiver-relative.
\item The \textbf{audit and self-model band} is where a frame that can represent
the relation attributes source, receiver, provenance, purpose, or intent. Its
facts are facts \emph{about} a run, held in an audit frame or a self-model, and
are never constituents of the run they describe.
\end{itemize}

\begin{definition}[Stratal facticity]\label{def:stratal-facticity}
A concept has \emph{stratal facticity} at the band that supplies the minimum
factual frame in which it is realised. A concept licenses only what that minimum
frame already contains, and does not license claims proper to a higher band. In
particular: a substrate concept is a fact wherever the rule runs or the
deformation occurs, and needs no frame relation; a frame-relation concept is a
fact only relative to a specified receiver, aperture, and grammar; an
audit concept is a fact only in a specified frame that holds the relevant relation
and performs the attribution. A concept realised in one band may be
\emph{referred to} from a higher band, but the reference does not relocate the
fact downward, and no audit attribution may be projected down into the
constituent frames of the run it describes.
\end{definition}

\paragraph{Reading the tables.} The three tables share the columns
\textsc{Concept}, \textsc{minimum factual frame} (the smallest arrangement in
which the concept is a fact), \textsc{what it licenses}, and \textsc{what it does
not license}. They are deliberately split by band rather than merged into one
long table, because the split is the claim: the last column of each band is where
level errors are refused. The prose after each table ties the rows to the running
cases without repeating them per row.

\begingroup\small
\begin{table}[t]\centering\small
\caption{Substrate band: rules run and deformations occur at this grain.}
\label{tab:band-substrate}
\begin{tabular}{@{}>{\raggedright\arraybackslash}p{0.15\textwidth}>{\raggedright\arraybackslash}p{0.26\textwidth}>{\raggedright\arraybackslash}p{0.26\textwidth}>{\raggedright\arraybackslash}p{0.24\textwidth}@{}}
\toprule
\textbf{Concept} & \textbf{Minimum factual frame} & \textbf{What it licenses} & \textbf{What it does not license}\\
\midrule
Rule & A frame's $\Gamma$ that runs at a grain. & That a transition can occur when conditions hold. & That the rule means, intends, or addresses anything.\\
Deformation & A realised change of graph state under a rule. & That something happened at this grain. & That it was received, understood, or aimed.\\
Distinction & A difference in state, structure, or admissible continuation. & That there is something to carry or hold. & That any frame has taken it in.\\
Grain & The specified resolution at which a frame runs. & Relativising every other concept to a resolution. & A privileged or grainless view.\\
Causal influence & A carried distinction realised as a transition in a receiving frame. & Calling the event communication at that grain. & A kind-split between ``mere cause'' and communication, or an intrinsic sender.\\
\bottomrule
\end{tabular}
\end{table}
\endgroup

The substrate band is where the sodium channel actually gates: a
conformational deformation follows a voltage distinction under the channel's rule,
and that is a fact whether or not any downstream frame reads it as a signal. It is
where the punch lands and the tap contacts---kinetic deformation at bodily
grain---and where a stored record is physical structure before it is a message. In
Ocelot's internal implementation, it is where a data structure is written or a
procedure runs; that a run occurred is a substrate fact, but (per
Remark~\ref{rem:ocelot-caveat}) it licenses no claim that the implementation
proves the substrate.

\begingroup\small
\begin{table}[t]\centering\small
\caption{Frame-relation band: a difference carried, exposed, or held across a boundary, relative to a receiver.}
\label{tab:band-frame}
\begin{tabular}{@{}>{\raggedright\arraybackslash}p{0.17\textwidth}>{\raggedright\arraybackslash}p{0.25\textwidth}>{\raggedright\arraybackslash}p{0.25\textwidth}>{\raggedright\arraybackslash}p{0.24\textwidth}@{}}
\toprule
\textbf{Concept} & \textbf{Minimum factual frame} & \textbf{What it licenses} & \textbf{What it does not license}\\
\midrule
Boundary / crossing & A relation between two frames whose condition is known by crossing. & That a difference can traverse. & That the crossing was intended or attributed to a source.\\
Carrier / trace / record & Realised transition or held structure exposing a distinction. & That a difference is available to be recovered. & That its origin or meaning travels with it.\\
Holding / exposure & Structure persisting or offered at a boundary. & Communication without uptake. & That any receiver has acted.\\
Aperture & Receiver-relative restriction on intake and action. & Bounding what can enter this receiver. & A claim about other receivers or a total intake.\\
Grammar / admissibility & The receiver's rules for what it can admit. & Deciding what is admissible at this grain. & Semantic truth or provenance.\\
Availability & A distinction the receiver can retain or discriminate. & Conditioning a later rule on the carrier. & That the distinction is operative or acted upon.\\
Uptake / operative reception / action & Availability that participates in a rule, state, or action. & That a difference did work in this receiver. & Crediting a receiver from an external description alone.\\
Encoding / decoding & A grammar mapping distinctions to and from a carrier form. & Recovery under a shared or installed scheme. & An intrinsic or universal code in the carrier.\\
Token / label / polarity & Carrier features a grammar reads as discrete values. & Cheap selection of a held shape. & That the label is self-interpreting.\\
Lexical edge & An installed selector to a held shape at a grain. & Fast retrieval when the edge is present. & That the shape equals the edge, or that absence erases the shape.\\
Meaning / content / interpretation & A receiver-side distinction made under a grammar. & Receiver-relative content at a grain. & A meaning intrinsic to the carrier or shared by fiat.\\
Error / residual / integrity & A mismatch or non-admissible remainder under a rule. & Detecting a scoped failure. & Global storage of all that was discarded, or source truth.\\
Repair / correction / feedback & A differential condition altering a later transition. & Discard, retransmit, correct, or degrade. & That two-way interaction is required for communication.\\
Common grammar / codebook & A grammar realised across both frames. & Coordinated recovery where it is held. & A free or universal precondition of communication.\\
Addressing / routing & Carrier features a relation uses to select a next frame. & Delivery toward a frame at a grain. & Proof that the addressee is the true or intended receiver.\\
Time / order & A partial order induced by a receiver's rules over crossings. & Sequencing relative to a receiver. & A total order given free by the medium.\\
\bottomrule
\end{tabular}
\end{table}
\endgroup

The frame-relation band is where the network link-layer frame lives: flags,
stuffing, headers, a frame-check sequence, addressing, buffers, and
acknowledgement state are all frame-relation facts, admissible or operative only
at their layer's grain (Section~\ref{sec:networking}). The tap and punch sit here
too once a receiver is specified: the same kinetic deformation is available under
one bodily or social grammar and merely felt under another. The sodium channel's
gating becomes a frame-relation fact the moment a downstream frame conditions on
the resulting current. A record is a held carrier whose reading is a later,
separate crossing. An Ocelot utterance and its lexicon belong here as a token or
lexical edge selecting a held shape under a reader's installed grammar; a word in
human language is the same relation at a learned grain. None of these rows names a
source or an intent: those are audit facts, tabulated next.

\begingroup\small
\begin{table}[t]\centering\small
\caption{Audit and self-model band: attributions about a run, held in a frame that represents the relation.}
\label{tab:band-audit}
\begin{tabular}{@{}>{\raggedright\arraybackslash}p{0.17\textwidth}>{\raggedright\arraybackslash}p{0.25\textwidth}>{\raggedright\arraybackslash}p{0.25\textwidth}>{\raggedright\arraybackslash}p{0.24\textwidth}@{}}
\toprule
\textbf{Concept} & \textbf{Minimum factual frame} & \textbf{What it licenses} & \textbf{What it does not license}\\
\midrule
Source / receiver attribution & An audit frame that directly perceives a crossing or reads an attachment. & Naming ends of a relation for a purpose. & That the run itself contains a source or receiver participant.\\
Direct perception & An audit frame that itself holds the crossing relation. & First-hand attribution of an end. & Attribution where the relation is not held.\\
Explicit attachment & A paid carrier asserting an origin (signature, header, label). & Evidence toward a source. & Proof of provenance; it can be forged or copied.\\
Implicit attachment & Context or convention a frame reads as origin evidence. & Weaker evidence toward a source. & Certainty, or a fact independent of the reading frame.\\
Provenance & An audit chain of attachments and perceptions. & Accumulating frame-relative evidence of origin. & Closure of ultimate source from a delivered carrier alone.\\
Identity / authentication & Rules testing an attachment against held keys or records. & A scoped predicate on claimed identity. & That the predicate settles truth or intent.\\
Sender / receiver & Roles an audit frame assigns to attributed ends. & Talking about the relation's asymmetry. & Intrinsic participants inside the run.\\
Purpose / function & An audit or selection ascription over a run. & Explaining or selecting for an effect. & An inner aim operative in the substrate step.\\
Intent & A self-model or audit inference about a run's disposition. & Predicting or judging a frame from its course. & An inner operator, or a precondition of communication.\\
Truth & An audit relation between a claim and a held relation. & Assessing a correspondence where the relation is held. & A view from nowhere or a carrier-intrinsic fact.\\
Audit frame & A specified frame that can represent the relations. & Making any attribution above. & A privileged observer (Axiom~\ref{GE-ax:no-view}).\\
Self-model & A frame's incomplete model of its own run. & Calibrated attribution of one's own intent or reception. & Completeness, or transparency to the run it models.\\
Social convention & A grammar and attachment scheme held across many frames. & Shared, cheaper attribution and codebooks. & Intrinsic meaning or a universal precondition.\\
\bottomrule
\end{tabular}
\end{table}
\endgroup

The audit band is where a network operator attributes a frame to a MAC address,
where a signature or header is read as an origin attachment, and where a
diagnostic log asserts provenance---each an audit fact scoped to a frame that
holds the relevant relation, never a constituent of the transmitted frame itself.
It is where a reader of a record attributes an author, where a hearer of a word
attributes a speaker's intent, and where an Ocelot audit or self-model would
attribute source, purpose, or intent to one of its own runs. Crucially, the ends
named ``sender'' and ``receiver'' are assigned here, not found in the substrate
step; the run proceeds whether or not any frame ever performs the attribution.

\begin{remark}[Attribution is not participation]\label{rem:attribution-not-participation}
Source, receiver, sender, message, purpose, and intent are not intrinsic
participants in a communication. They are attributions made in an audit band.
An audit frame may attribute a source or a receiver to a crossing only if it
either \emph{directly perceives} that crossing---by itself holding the relevant
relation---or reads an \emph{explicit} or \emph{implicit attachment} carried with,
or contextually bound to, the difference. Absent direct perception or a readable
attachment, a frame has evidence local to one side and no licensed attribution of
the other. This is the relativity of grounding (Axiom~\ref{GE-ax:no-view})
applied to origin and destination: naming the ends of a relation records an
asymmetry an audit frame draws, and does not insert a participant into the run.
Consequently, the roles of the minimal relation (Section~\ref{sec:minimal-relation})
need not be represented by any constituent frame of the run; a substrate step can
be fully a communication at its grain while no frame within it holds a
representation of ``source'' or ``receiver'' at all.
\end{remark}

\begin{remark}[Attachments are paid evidence, not proof]\label{rem:attachments-evidence}
An attachment---a signature, a header field, a return address, a byline, a
contextual cue---is itself a carrier that was paid for at some crossing. It is
therefore \emph{evidence} toward a source or receiver, weighted by the cost of
producing and of forging it, and not a \emph{proof} of provenance. Explicit
attachments can be copied, replayed, or fabricated; implicit attachments depend on
the reading frame's conventions and can mislead. Authentication rules raise the
cost of a false attachment but still yield a scoped predicate, not closure of
ultimate origin (cf.\ the provenance limit, Proposition~\ref{prop:network-layered}
and Section~\ref{sec:repair}). Accumulating attachments accumulates frame-relative
evidence; it never manufactures a view from nowhere.
\end{remark}

\begin{remark}[Intent is a self-model inference about a run]\label{rem:intent-selfmodel}
Intent is not an inner operator in the substrate step and is not a condition on
communication. It is an inference---made in an audit frame or, reflexively, in a
frame's own self-model---about the \emph{disposition} a run exhibits: what the run
does, and would do, across the conditions the modeller can represent. This applies
the mind-and-self grounding's self-model core: a self-model is incomplete by
construction (\citealp[Prop.~\ref{MS-prop:selfmodel}, Sec.~\ref{MS-sec:selfmodel}]{}),
so an attribution of one's own intent is a calibrated wager, not a readout of a
hidden aim; and there is no dedicated primitive for influence beneath relation
(\citealp[Prop.~\ref{MS-prop:one-influence}]{}), so intent adds no operator to the
run. The constraint the fable of the reflexive animal makes vivid holds here: a
run's nature is read from what it does under its conditions, not from a declared
inner reason, and confining or reframing the conditions does not convert an
attributed intent into a substrate participant. An attributed intent may inform
prediction and judgement in the audit band; it may never be projected down into
the constituent frames as though it were a rule they ran.
\end{remark}

\begin{remark}[Ocelot as an internal implementation example]\label{rem:ocelot-caveat}
Where this paper mentions Ocelot's utterance pipeline, lexicon, records, or audit,
it does so as a clearly labelled \emph{internal implementation example}, not as
independent evidence for the substrate. That an implementation writes a record,
runs a lexical lookup, or emits an utterance is a substrate and frame-relation
fact about that system; it does not license the claim that the implementation
\emph{proves} the substrate, and any attribution of source, purpose, or intent to
one of its runs remains an audit-band inference subject to
Remark~\ref{rem:attribution-not-participation} and
Remark~\ref{rem:intent-selfmodel}. The utterance pipeline is therefore used only
to instantiate the token/lexical-edge and record rows of
Table~\ref{tab:band-frame}; source attribution over it is not asserted here,
because doing so would overclaim.
\end{remark}

\subsection{Worked case: a frozen language model, a prompt-installed term, and a persistent update}\label{sec:llm-worked-case}

The stratigraphy is easiest to misread where a token is treated as a property an
input carries with it. This worked case fixes the reading against a standard,
publicly documented system---a \emph{frozen} large language model---and does so
with the strict graph-law definition of tokenisation stated first. Nothing below
is offered as novel to language-model practitioners, and nothing below is offered
as proof of the substrate; the case only exhibits the token/lexical-edge rows of
Table~\ref{tab:band-frame} and the stratal-facticity discipline of
Definition~\ref{def:stratal-facticity} in a familiar mechanism.

\paragraph{The strict definition.} Tokenisation at a frame $F$ is the matching of
an input made available at $F$'s boundary and grain to an \emph{already extant,
admissible edge held by $F$}. A token is therefore the frame-relative successful
match, not an intrinsic property of the input. Three consequences follow, and
they must be kept apart because they occupy different rows of the band.
\begin{itemize}[leftmargin=*]
\item An input that $F$'s aperture excludes is \emph{absent} for $F$: there is no
carrier at $F$'s boundary to match, so no question of tokenisation arises.
\item An input that reaches the relevant grain but finds \emph{no matching
admissible edge} is not tokenised there. It may be collapsed to a coarse unknown
category, left as residual (Definition~\ref{GE-def:residual}), or trigger a
separately paid extension of the grammar. None of these is tokenisation.
\item Minting a new edge is not cheap tokenisation. Installing a durable new
selector is a distinct write charged where it occurs
(Proposition~\ref{prop:pays-at-crossing}); tokenisation to an existing edge is
the cheap selection the token row licenses, and the two must not be conflated.
\end{itemize}

\paragraph{Frame~1: the fixed tokenizer and model.} A deployed frozen model holds
a fixed tokenizer vocabulary and fixed parameters. Its tokeniser can match an
input only to pre-existing vocabulary or token-ID edges; ordinary inference mints
no persistent vocabulary edge and no persistent parameter edge. This is exactly
the fixed-vocabulary constraint that fixed-vocabulary neural language models make
explicit \citep{bengio2003}. Two situations that look alike must be separated at
this frame. A genuinely \emph{novel token} would require an edge the frame does
not hold, and the frozen frame cannot supply it during inference. A \emph{novel
string}, by contrast, is routinely representable as a novel \emph{sequence of
existing} subword or byte units \citep{sennrich2016,kudo2018,radford2019}: the
string is new, but every edge used to segment it is already held, so this is
successful tokenisation to extant edges, not minting. The distinction between a
new string over old units and a new unit is the whole of the point.

\paragraph{Grain, not exemption from the rule.} Early closed word-vocabulary
systems make the rule visible by failing loudly: an out-of-vocabulary word form
finds no matching admissible edge at the \emph{word} grain and is collapsed to a
coarse unknown category, the \texttt{UNK} token \citep{luong2015}. Modern subword
and byte-level systems do not repeal this rule; they \emph{lower the grain} at
which matching succeeds. Byte-level schemes push the grain down to a small,
closed set of byte edges so that essentially any input segments to held units
\citep{radford2019}, and subword schemes sit between word and byte grains
\citep{sennrich2016,kudo2018}. The coarse-unknown outcome has become rare in
practice not because the graph law changed but because the admissible edge set was
made fine enough that a match almost always exists. The residual and coarse-unknown
possibilities of the strict definition remain the correct account of what happens
when they do not.

\paragraph{Frame~2: the runtime context.} A prompt can install a temporary
association---``let \emph{glorp} denote the following shape''---between an existing
token sequence and a prompt-provided definition, and the model can condition its
continuation on that association without any gradient update
\citep{brown2020}. In stratal terms this is a frame-relation fact local to the
runtime context frame: a receiver-side availability (Definition~\ref{def:availability})
held for the duration of the context. It is real, and it is paid for---in context
capacity and in the conditioning it induces---but it changes neither the tokenizer
vocabulary nor the model weights. The token sequence for \emph{glorp} is still
matched to the same pre-existing subword edges as in Frame~1; what is new is a
context-local relation between that sequence and a shape, not a new lexical edge in
the fixed frame. This is precisely the paper's distinction between a context-local
availability and a persistent update: the term is operative in this run and absent
from the next fresh context.

\paragraph{Frame~3: the containing persistent-update frame.} A durable new relation
can be paid for and held, but only in a larger frame than ordinary inference. Three
such frames are standard. A tokenizer/vocabulary change with a matching embedding
update mints a genuinely new token edge and pays for its parameters. Training or
fine-tuning revises parameter edges so that a relation becomes held in the weights.
An external memory or retrieval write records a durable relation in a datastore the
model reads at inference, so that a nearest-neighbour or retrieval step can recover
it later \citep{khandelwal2020}. Each is a minting or update write in the sense of
Proposition~\ref{prop:pays-at-crossing}, charged in the containing frame; none is
performed by the frozen inference frame on its own. The context installation of
Frame~2 and the persistent update of Frame~3 are the same distinction the paper
draws throughout between an availability held for a crossing and a relation held
across successions.

\paragraph{Caveats.} Some deployed systems do perform online learning, periodic
fine-tuning, or external-memory writes during operation
\citep{khandelwal2020}; those are larger-frame operations in the sense above, not
counterexamples to the account of ordinary frozen inference. The mapping claims
only what the stratigraphy already says: tokenisation is a frame-relative match to
a held admissible edge; a coarse-unknown category, a residual, or a paid grammar
extension are the distinct outcomes when no such edge exists; and a context-local
association is not a persistent lexical or parameter edge. These are facts about
the model as a frame relation and its containing update frames, and they are used
here only to instantiate the token and lexical-edge rows of
Table~\ref{tab:band-frame}, not to assert any implementation-specific claim about
Ocelot or to prove the substrate.


\section{Network frames: a literal engineering grounding}\label{sec:networking}

Networking is not a decorative analogy for this paper. At the data-link layer,
engineers literally call the bounded unit of transmission a \emph{frame}. The
standards supply a mature, independently developed case in which boundary,
grammar, admissibility, opacity, residual, finite aperture, and layered
containment are operational rather than metaphorical.

Terminology must be exact. RFC~1662 calls a \emph{frame} the data-link-layer unit
of transmission; a \emph{datagram} is a network-layer unit such as an IP
datagram; and a \emph{packet} is ordinarily the unit of encapsulation passed
between the network and data-link layers \citep{simpson1994}. A frame is therefore
not a universal name for any bundle of bits. The same bounded configuration can
be a link-layer frame under one grammar and an opaque payload carrying a
higher-layer packet or datagram under another. This is a direct engineering
instance of the distinction between a local frame and a containing one.

Ethernet MAC framing supplies destination and source addresses, a length or
EtherType field, payload and padding, and a frame check sequence; the standard
specifies the MAC frame separately from physical-layer preamble and start-frame
delimiter \citep{ieee8023}. PPP supplies an especially transparent grammar: flags
mark a boundary, byte or bit stuffing prevents payload from forging that boundary,
and a frame-check sequence detects specified corruptions \citep{simpson1994}.
These mechanisms do not make a frame meaningful in a human sense. They make a
boundary decidable and integrity testable at a particular link-layer grain.

\begin{proposition}[Network admissibility is layer-relative]\label{prop:network-layered}
A link-layer frame may be admissible under its framing and integrity grammar while
its payload remains opaque, malformed, or unusable under a receiver's next-layer
grammar. Conversely, an apparent higher-layer packet cannot be received as such
without some lower-layer framing relation making a carrier available. Integrity,
addressing, and interpretation are therefore scoped to the layers that realise
them.
\end{proposition}

The proposition is an arrangement claim about a known engineering fact, not a
claim that Ethernet proves the substrate. The distinction is visible in the
checks themselves. Ethernet's frame check sequence covers the MAC frame; IPv4's
header checksum covers the IP header rather than the whole payload
\citep{ieee8023,postel1981}. A passed check says that a scoped predicate was met.
It does not establish source truth, application semantics, or even a valid next
layer. An unknown EtherType or protocol number can leave payload present but
non-operative at a receiver's current grammar.

An invalid PPP frame is equally instructive. It can be discarded by the PPP
receiver because it violates its frame grammar or check sequence. That does not
make it an absence of communication in the broad sense of this paper. Physical
differences have crossed; the failed frame may change an error counter or trigger
an interface rule; a diagnostic frame may retain a residual. That residual is
per-rule---it is what a particular receiver rule could not admit at its grain
(Definition~\ref{GE-def:residual})---not a global store of everything discarded.
Any retained diagnostic form is a separately paid later record or copy in some
holder, charged where that retention occurs, rather than a standing residual the
relation keeps for free. It is instead
non-admissible or non-operative \emph{as a PPP payload carrier} at that grain. The
same discipline prevents a reader from treating a human message ignored by its
addressee as though no difference had been communicated at all.

TCP makes the payment structure visible at a larger layer. Sequence numbers,
acknowledgements, retransmission timers, receive windows, and sender/receiver
buffers retain distinctions across successive crossings \citep{eddy2022,paxson2011}.
They pay for a particular form of reliable, ordered reception. The ordering they
provide is a partial order carried by the receiver's rules over its crossings,
not a total order given free by the medium (Proposition~\ref{GE-prop:partial-order}).
They are not
necessary to communication as such: an IP datagram, a UDP datagram, a one-way
broadcast, or a forward-error-corrected transmission can make a difference
available without a return acknowledgement or a durable sender-side copy. The
end-to-end argument makes the limit principled. Correctness that depends on
application knowledge cannot be supplied conclusively by a lower layer, even
when lower layers provide useful checks and redundancy \citep{saltzer1984}.

This gives four constraints that the general account must respect. First,
framing grammar is not semantic truth. Second, a receiver's acceptance of a
carrier is not a proof of provenance. Third, repair can be layered and
contingent, not a universal prerequisite. Fourth, no frame may be credited with
higher-layer distinctions merely because a containing network observer can read
the bits. The name \emph{frame} in networking is not evidence for Graphic
Equalisation; the operational recurrence is evidence only for the limited
composition claim that a common shape is doing work across domains.

\section{Records, code, signalling, and lexical edges}\label{sec:instances}

The network case is deliberately first because it is precise, non-human, and
literal. It is not the sole case. The following instances are not offered as
separate origins of a communication primitive. They are ways the same inherited
frame machinery is independently made concrete.

\paragraph{Records and code.} A record carries a difference across a separation in
successions. Inscription pays where the holder writes a trace; preservation pays
where the holder retains it; reading pays where a later frame constructs a
receiver-side distinction. Source code provides a sharper case. The source file,
parser, abstract syntax tree, compiled artifact, instruction stream, and running
machine are nested carriers at distinct grammars. A token may be irrelevant to a
processor, decisive to a parser, and decomposable in a compiler's intermediate
representation. That does not make the token decorative: at its own grain it is a
cheap edge to a held syntactic or semantic shape. The code case also denies a
common error. An executable file can communicate a reusable capability even when
it is not currently run; running it is a later, separately paid transition.

\paragraph{Molecular and cellular signalling.} Ligand binding, ion gradients,
hormone reception, gene regulation, and quorum-sensing systems make clear that a
receiver grammar need not be linguistic or conscious. The relevant grammar may
be receptor geometry, concentration sensitivity, timing, or a downstream
regulatory rule. A molecular event is kinetic communication before it is a
biological signal in any special historical sense. Evolutionary signalling theory
can then ask an additional question---how sender and receiver dispositions became
stabilised---without defining the cross-frame communication relation itself
\citep{maynardsmithharper2003,scottphillips2008}.

\paragraph{Animal signals and social convention.} Alarm calls, displays, scent
marks, and contact cues are richly useful tests because they combine a physical
carrier, receiver-side discrimination, interaction, deception, and sometimes
shared interest. Their biological history may distinguish a signal from a cue in
the specialist literature. That distinction remains important for explaining
reliability and selection, but it does not licence the paper to call a predator's
impact on prey non-communicative merely because it is coercive, or to require an
intentional sender before a receiver-side difference can exist.

\paragraph{Human language.} Language is a high-density, learned case in which a
small token can make a large held shape available under a reader's lexical and
grammatical practices. Its significance here is not that language creates
communication. It makes the costs of naming, recovery, ambiguity, repair, and
cross-domain coordination unusually inspectable. A project term is operative
compression at a reader's lexical grain. It must be installed in prose,
definition, mechanism, and boundary; a lexical edge is not a proof, and it is not
a substitute for the shape it names.

\subsection{Exact retrieval and constrained reconstruction}\label{sec:retrieval}

An exact token match is cheap when the receiver has that token in an installed
grammar. The receiver need not search a universal space: the token selects an
edge already held at its grain. When exact retrieval fails, recovery need not
become global reconstruction. A held target shape can constrain a local search
among nearby candidate edges, partial forms, related contrasts, or routes through
a grammar. This is a structural possibility, not an unsupported claim that every
system has one particular search algorithm or complexity class.

Tip-of-the-tongue is the bounded human limit case. A speaker can hold a target's
meaning and portions of its phonology or syntax while failing to fetch the full
lexical form \citep{brownmcneill1966,levelt1999}. The phenomenon is not itself a
cross-frame event unless one chooses an internal subsystem boundary; it is better
used as evidence that receiver-side availability is layered. The lexical edge can
be unavailable while the held shape remains sufficiently present to guide a
close-enough recovery. Its value is precisely that it does not require a thesis
about consciousness or qualia.

\begin{remark}[Carrier identity is not success]\label{rem:identity}
Exact preservation of a carrier is one mode of communication, not its definition.
A packet can be re-encoded across a different physical link; a program can be
compiled; a record can be copied; an explanation can be paraphrased; a receptor
can act on a concentration range rather than one molecule's identity. What is
preserved is assessed only at a specified grain and in a specified audit frame.
No such assessment licenses a claim of absolute message identity.
\end{remark}

\section{Residual, repair, and the limits of reception}\label{sec:repair}

Residual is where the account becomes more than a vocabulary of carriers. A
receiver can fail to admit a difference, can admit it without recovering a later
shape, can recover a shape without acting, or can act on a corrupted or partial
reconstruction. These are not one failure mode. They occur at different grammars,
and their distinction is what stops a system's external description from being
silently charged to its internal reception.

Error correction supplies a clear family. A parity check, checksum, syndrome,
missing acknowledgement, or violated prediction creates a differential condition
under a receiver grammar. A later rule may discard, correct, request another
carrier, select an alternative reconstruction, or continue with degraded
operation. Most such repair runs under an existing grammar: it fires rules the
receiver already holds, and pays only for the comparison, correction, or
retransmission those rules specify. Extending the grammar itself---installing a
new rule, decoder, or lexical edge so that a previously inadmissible difference
becomes recoverable---is a distinct and separately paid write, and is only one
repair mode among the others. The paper's claim is deliberately weak enough to
include both
feedback and forward correction: successful correction requires some mismatch or
constraint to alter a later state transition; it does not require that mismatch to
be conscious, visible to an auditor, returned to the source, or held indefinitely.

Conversation repair and protocol repair are therefore analogous without being
identical. In conversation, a recipient can request clarification, offer a
candidate, or correct a reference \citep{clark1996,clarkwilkesgibbs1986}. In a
network, a receiver can reject a malformed frame, acknowledge a sequence number,
or correct a decodable codeword. Both expose a residual to a later rule. Neither
supports the stronger claim that all communication requires two-way interaction:
a written record, broadcast, or molecular gradient may be received without a
reply, and may be carried without reception at all.

The same constraint limits provenance. A receiver can often establish that a
carrier meets its own integrity predicate. It cannot thereby settle, from inside,
why the carrier was produced, whether every relevant source condition was known,
or whether an external description's intended interpretation is correct. This is
the general prohibition on a frame closing its own provenance, applied at the
communication boundary. Authentication, signatures, logs, and social testimony
are additional carriers and rules; they can improve a frame's local evidence but
do not create a view from nowhere.

\section{Prior art, contribution, and exclusions}\label{sec:prior-art}

The paper does not claim to discover transmission, signalling, coding, or
language. Its contribution, if the construction survives scrutiny, is an
arrangement: the same frame-relative economy describes kinetic influence,
records, protocol frames, and lexical recovery without treating one human case as
the foundation of the rest.

Shannon's mathematical theory gives an indispensable account of encoding,
channels, noise, capacity, and reliable transmission, while explicitly bracketing
semantic aspects of communication \citep{shannon1948}. Rate--distortion theory
shows why recovery can be graded and task-relative rather than exact
\citep{coverthomas2006}. Neither supplies a frame-relative account of what a
receiver can hold, what an audit frame may claim about provenance, or where a
particular cross-frame transformation pays. Communication complexity instead
measures what distributed parties must exchange to compute a function of their
separate inputs \citep{yao1979,kushilevitznisan1997}; it is a useful lower-bound
literature for specific cooperative tasks, not a definition of all communication.

Semiotic and pragmatic traditions correctly deny that a sign is self-contained.
Peirce's interpretant and Millikan's consumer-side account are especially useful
warnings against treating a token as sufficient by itself \citep{peirce1931,millikan1984}.
Gricean intention and later ostensive accounts illuminate a demanding human case,
but they must not become a universal gate that excludes punches, protocols, or
cellular signalling. Conversely, an account broad enough to include all causal
influence must not pretend that all causal differences possess the same
conventionality, reliability, or reconstructive richness. Those are contingent
properties of particular grammars and histories.

Networking prior art is direct rather than metaphorical. The standards and RFCs
already specify frame boundaries, transparent encoding, layered encapsulation,
integrity tests, retransmission, finite receive state, and endpoint limits
\citep{ieee8023,simpson1994,postel1981,eddy2022,saltzer1984}. This paper does not
improve those specifications. It uses them as a stringent independent case: if
the proposed grounding cannot distinguish link-layer admissibility from
application-level correctness, or cannot account for a frame becoming opaque
payload at another layer, its vocabulary is not doing the work claimed.

Consciousness, phenomenal character, and qualia are out of scope, as they are in
the mind-and-self grounding. So is a claim that communication requires a
conceptual space, a global semantic geometry, or the free-energy principle.
Those may be useful domain models under additional commitments; none is required
to state a difference crossing and paying at a frame boundary.

\section{Falsifiability commitments and evidence}\label{sec:falsifiability}

The framework-level identification of causal influence with communication is not
presented as an ordinary laboratory hypothesis. It is a way of refusing an
otherwise unmotivated kind split: a causal influence is a carried distinction at
some grain. The claims below are narrower promoted commitments. A failure bounds
this grounding's range or replaces one mechanism; it does not prove or disprove
the substrate from a view outside every frame.

\begin{enumerate}[leftmargin=*]
\item \textsc{[shape]} \textbf{Payment is located at a frame change.} The
application fails if a realised receiver-side distinction can be supplied with no
write in the receiver and no sustaining or coupling write in the containing
relation, after the relevant frames are specified. A later observer's compact
description is not a counterexample.
\item \textsc{[shape]} \textbf{A causal influence is communicable difference at a stated grain.}
The identification is a structural vocabulary unification, not a claim that every grain supplies the
same recovery or semantic description. For a named receiving frame, boundary, and grain, the
construction fails if a realised causal influence requires an element not representable as a carried or
changed relation at that boundary, and no finer stated grain renders it a cross-frame difference. An
intra-frame change with no carried distinction across the selected boundary is excluded from the
communication relation at that grain; it may be communication only under a different, explicitly named
frame relation. Mere unavailability of recovery at the chosen grain is not a counterexample.
\item \textsc{[detail]} \textbf{Admissibility is layer-relative.} The network
mapping fails if a single validity predicate is found to be authoritative across
all layers of an actual stack, so that no carrier can be accepted at one grammar
while opaque, malformed, or inoperative at another.
\item \textsc{[detail]} \textbf{Recovery is graded and can outrun exact token
identity.} The account's retrieval mechanism fails if reliable receiver action
requires exact carrier identity in every relevant case, or if target-shape access
never constrains a local lexical or structural recovery when direct access fails.
\item \textsc{[detail]} \textbf{Repair uses a differential condition.} The
repair account fails if a system corrects a carrier-dependent error while no
mismatch, syndrome, constraint violation, or alternative comparison affects any
later state transition. Feedback need not occur, so a forward-error-correcting
code is not a counterexample.
\item \textsc{[shape]} \textbf{Provenance is not closed by delivery.} The
account fails if a receiver can settle from its own delivered carrier alone the
ultimate source, truth, or complete intended interpretation of that carrier,
rather than accumulating only additional frame-relative evidence. Attachments
(signatures, headers, return addresses, contextual cues) are paid evidence toward
a source, not proof of it (Remark~\ref{rem:attachments-evidence}); the account
also fails if any single attachment is shown to be, by itself, incorrigible
proof of provenance rather than forgeable, replayable, or convention-dependent
evidence.
\item \textsc{[shape]} \textbf{Audit attributions do not project into
constituents.} The stratigraphy fails if a licensed communication requires that
one of its constituent frames represent a source, sender, receiver, message,
purpose, or intent---that is, if some causal influence cannot be a communication
at its grain unless an audit-band attribution is already realised as a rule those
frames run (Definition~\ref{def:stratal-facticity},
Remark~\ref{rem:attribution-not-participation}). Attributing intent, source, or
purpose from an audit frame or self-model is not a counterexample; projecting one
of those attributions downward into the run as an operative participant would be.
\end{enumerate}

The evidence plan has corresponding layers. First, a formal audit should show
that every term in Definition~\ref{def:communication} is realised through
inherited graph, rule, aperture, residual, and boundary structure, without a new
primitive. Second, executable protocol fixtures can demonstrate framing,
encapsulation, corruption, discard, retransmission, and layer-specific parsing;
they establish specification and implementation adequacy, not the substrate's
truth. Third, standards and independently developed biological, computational,
and linguistic cases can test whether the arrangement transfers. Finally,
lexical-access experiments can test the narrow claim about partial recovery. None
of these evidences should be silently promoted from a case of composition to a
proof of the framework.

\section{Conclusion}\label{sec:conclusion}

Communication is difference across frames. It need not begin with a human mind,
an intention, a word, a signal convention, or a recipient's success. A tap on the
shoulder and a punch establish the ordinary case: kinetic influence is
communication, differing from language in magnitude, grain, persistence, and
recoverable structure rather than in kind. A record can communicate while unread;
a network frame can be received while semantically opaque; a word can be held as
a shape while its lexical edge cannot be fetched.

What unifies these cases is not a free carrier or a view from nowhere. Each
operative change pays where it changes frames. A boundary crossing, a receiver's
admission, a retained record, and a repair operation are separately chargeable
writes. The carrier is not its own meaning, and a higher frame's description does
not make distinctions free to the lower one. Networking makes this especially
clear because its frames are literal, bounded, layered, and independently
specified: a link-layer frame, a network-layer datagram, and an application
message are not interchangeable even when they travel in one physical event.

The proposed grounding therefore earns its place only if it lets the same small
substrate vocabulary do work that the local literatures leave distributed: locate
the grain of a communication, the frame that pays for it, the grammar that can
receive it, the residual that can repair it, and the boundary beyond which
provenance cannot be closed. If it cannot make those distinctions without adding
a communication primitive, treating causal influence as a different kind, or
silently importing an external interpreter, the draft should be narrowed or
rejected.

\begin{thebibliography}{99}

\bibitem[Bateson(1972)]{bateson1972}
Gregory Bateson.
\newblock \emph{Steps to an Ecology of Mind}.
\newblock Chandler Publishing Company, San Francisco, 1972.

\bibitem[Bengio et~al.(2003)]{bengio2003}
Yoshua Bengio, R\'ejean Ducharme, Pascal Vincent, and Christian Jauvin.
\newblock A neural probabilistic language model.
\newblock \emph{Journal of Machine Learning Research}, 3:1137--1155, 2003.

\bibitem[Brown and McNeill(1966)]{brownmcneill1966}
Roger Brown and David McNeill.
\newblock The ``tip of the tongue'' phenomenon.
\newblock \emph{Journal of Verbal Learning and Verbal Behavior}, 5:325--337, 1966.

\bibitem[Brown et~al.(2020)]{brown2020}
Tom B. Brown, Benjamin Mann, Nick Ryder, Melanie Subbiah, Jared Kaplan, et~al.
\newblock Language models are few-shot learners.
\newblock In \emph{Advances in Neural Information Processing Systems~33 (NeurIPS 2020)}, pages 1877--1901, 2020.

\bibitem[Clark(1996)]{clark1996}
Herbert H. Clark.
\newblock \emph{Using Language}.
\newblock Cambridge University Press, Cambridge, 1996.

\bibitem[Clark and Wilkes-Gibbs(1986)]{clarkwilkesgibbs1986}
Herbert H. Clark and Deanna Wilkes-Gibbs.
\newblock Referring as a collaborative process.
\newblock \emph{Cognition}, 22(1):1--39, 1986.

\bibitem[Cover and Thomas(2006)]{coverthomas2006}
Thomas M. Cover and Joy A. Thomas.
\newblock \emph{Elements of Information Theory}.
\newblock Second edition, Wiley, Hoboken, 2006.

\bibitem[Eddy(2022)]{eddy2022}
Wesley Eddy, editor.
\newblock \emph{Transmission Control Protocol}.
\newblock RFC 9293, Internet Engineering Task Force, August 2022.

\bibitem[IEEE(2022)]{ieee8023}
Institute of Electrical and Electronics Engineers.
\newblock \emph{IEEE Standard for Ethernet}.
\newblock IEEE Std 802.3-2022, New York, 2022.

\bibitem[Khandelwal et~al.(2020)]{khandelwal2020}
Urvashi Khandelwal, Omer Levy, Dan Jurafsky, Luke Zettlemoyer, and Mike Lewis.
\newblock Generalization through memorization: Nearest neighbor language models.
\newblock In \emph{International Conference on Learning Representations (ICLR 2020)}, 2020.

\bibitem[Kudo and Richardson(2018)]{kudo2018}
Taku Kudo and John Richardson.
\newblock SentencePiece: A simple and language independent subword tokenizer and detokenizer for neural text processing.
\newblock In \emph{Proceedings of the 2018 Conference on Empirical Methods in Natural Language Processing: System Demonstrations}, pages 66--71, 2018.

\bibitem[Kushilevitz and Nisan(1997)]{kushilevitznisan1997}
Eyal Kushilevitz and Noam Nisan.
\newblock \emph{Communication Complexity}.
\newblock Cambridge University Press, Cambridge, 1997.

\bibitem[Levelt et~al.(1999)]{levelt1999}
Willem J. M. Levelt, Ardi Roelofs, and Antje S. Meyer.
\newblock A theory of lexical access in speech production.
\newblock \emph{Behavioral and Brain Sciences}, 22(1):1--75, 1999.

\bibitem[Luong et~al.(2015)]{luong2015}
Minh-Thang Luong, Ilya Sutskever, Quoc V. Le, Oriol Vinyals, and Wojciech Zaremba.
\newblock Addressing the rare word problem in neural machine translation.
\newblock In \emph{Proceedings of the 53rd Annual Meeting of the Association for Computational Linguistics}, pages 11--19, 2015.

\bibitem[Maynard Smith and Harper(2003)]{maynardsmithharper2003}
John Maynard Smith and David Harper.
\newblock \emph{Animal Signals}.
\newblock Oxford University Press, Oxford, 2003.

\bibitem[Millikan(1984)]{millikan1984}
Ruth Garrett Millikan.
\newblock \emph{Language, Thought, and Other Biological Categories}.
\newblock MIT Press, Cambridge, Massachusetts, 1984.

\bibitem[Paxson et~al.(2011)]{paxson2011}
Vern Paxson, Mark Allman, Jerry Chu, and Matt Sargent.
\newblock \emph{Computing TCP's Retransmission Timer}.
\newblock RFC 6298, Internet Engineering Task Force, June 2011.

\bibitem[Peirce(1931)]{peirce1931}
Charles Sanders Peirce.
\newblock \emph{Collected Papers of Charles Sanders Peirce, Volume~2: Elements of Logic}.
\newblock Harvard University Press, Cambridge, Massachusetts, 1931.

\bibitem[Postel(1981)]{postel1981}
Jon Postel, editor.
\newblock \emph{Internet Protocol}.
\newblock RFC 791, Internet Engineering Task Force, September 1981.

\bibitem[Radford et~al.(2019)]{radford2019}
Alec Radford, Jeffrey Wu, Rewon Child, David Luan, Dario Amodei, and Ilya Sutskever.
\newblock Language models are unsupervised multitask learners.
\newblock Technical report, OpenAI, 2019.

\bibitem[Saltzer et~al.(1984)]{saltzer1984}
Jerome H. Saltzer, David P. Reed, and David D. Clark.
\newblock End-to-end arguments in system design.
\newblock \emph{ACM Transactions on Computer Systems}, 2(4):277--288, 1984.

\bibitem[Scott-Phillips(2008)]{scottphillips2008}
Thomas C. Scott-Phillips.
\newblock Defining biological communication.
\newblock \emph{Journal of Evolutionary Biology}, 21(2):387--395, 2008.

\bibitem[Sennrich et~al.(2016)]{sennrich2016}
Rico Sennrich, Barry Haddow, and Alexandra Birch.
\newblock Neural machine translation of rare words with subword units.
\newblock In \emph{Proceedings of the 54th Annual Meeting of the Association for Computational Linguistics}, pages 1715--1725, 2016.

\bibitem[Shannon(1948)]{shannon1948}
Claude E. Shannon.
\newblock A mathematical theory of communication.
\newblock \emph{Bell System Technical Journal}, 27:379--423 and 623--656, 1948.

\bibitem[Simpson(1994)]{simpson1994}
William A. Simpson, editor.
\newblock \emph{PPP in HDLC-like Framing}.
\newblock RFC 1662, Internet Engineering Task Force, July 1994.

\bibitem[Yao(1979)]{yao1979}
Andrew C.-C. Yao.
\newblock Some complexity questions related to distributive computing.
\newblock In \emph{Proceedings of the Eleventh Annual ACM Symposium on Theory of Computing}, pages 209--213. ACM, 1979.

\end{thebibliography}
\end{document}
